%!TEX root = ExerciseSheet2.tex
% -----------------------------------------------------------------------------
% PART A
% -----------------------------------------------------------------------------
\subsection*{Part A: The Idea}
\textit{These questions test the mathematical definitions and calculations devoid of physical context.}

\begin{enumerate}[label=\textbf{Q\arabic*.}, leftmargin=*]

    \item \textbf{Basic Probability \& Set Operations} \\
    Let $A$ and $B$ be two events in a sample space $S$. You are given the following probabilities:
    $P(A) = 0.40$, $P(B) = 0.35$, and $P(A \cap B) = 0.15$.
    \begin{enumerate}[label=\arabic*.]
        \item Calculate the probability of $A$ or $B$ occurring, $P(A \cup B)$.
        \item Calculate the conditional probability of $B$ occurring given that $A$ has occurred, $P(B|A)$.
        \item Are events $A$ and $B$ mutually exclusive? Justify your answer using the provided numbers.
        \item Are events $A$ and $B$ independent? Justify your answer mathematically.
    \end{enumerate}
    \vspace{0.3cm}

    \item \textbf{Conditional Probability \& Trees} \\
    An urn contains $5$ white balls and $3$ black balls. You draw two balls from the urn \textbf{without replacement}. Let $W_1$ be the event that the first ball is white, and $B_2$ be the event that the second ball is black.
    \begin{enumerate}[label=\arabic*.]
        \item Calculate $P(W_1)$.
        \item Calculate $P(B_2 | W_1)$. 
        \item Calculate the probability that the first ball is white \textbf{and} the second is black, $P(W_1 \cap B_2)$.
        \item Calculate the unconditional probability that the second ball is black, $P(B_2)$, by considering both possible outcomes for the first draw. \textit{(Hint: A tree diagram will help here).}
    \end{enumerate}

\end{enumerate}

\vspace{0.8cm}

% -----------------------------------------------------------------------------
% PART B
% -----------------------------------------------------------------------------
\subsection*{Part B: Exam Style Questions: Engineering Application}
\textit{These questions test the exact same mathematical concepts as Part A, but applied to the Electrical Engineering scenarios discussed in the lecture.}

\begin{enumerate}[label=\textbf{Q\arabic*.}, leftmargin=*, start=3]

    \item \textbf{The Addition Law \& Manufacturing Defects} \\
    A batch of freshly manufactured printed circuit boards (PCBs) is undergoing visual inspection. The probability that a board has a soldering defect is $0.08$. The probability it has a broken copper trace is $0.04$. The probability that a board has \textbf{both} defects is $0.01$.
    \begin{enumerate}[label=\arabic*.]
        \item What is the probability that a board has \textit{at least one} of these defects?
        \item If you randomly select a board and notice it has a broken copper trace, what is the probability that it also has a soldering defect?
    \end{enumerate}
    \vspace{0.3cm}

    \item \textbf{The Law of Independence \& System Redundancy} \\
    A critical data server requires a cooling fan to prevent thermal shutdown. A standard fan has a $5\%$ chance of failing during a year of continuous operation ($P(F) = 0.05$). To improve reliability, an engineer installs two identical fans in parallel. They operate completely independently of one another.
    \begin{enumerate}[label=\arabic*.]
        \item Using the Multiplication Law for independent events, calculate the probability that \textbf{both} fans fail during the year.
        \item The server will survive as long as \textit{at least one} fan is working. Using the Complement rule, what is the probability that the server survives the year without thermal shutdown?
        \item Based on these numbers, explain why EEEs heavily rely on the mathematics of independent events when designing critical systems.
    \end{enumerate}
    \vspace{0.3cm}

    \item \textbf{The Problem of \lq Luck\rq\ Resolved (Bridging the Gap)} \\
    In Module 1, we asked: \textit{"You flip a fair coin 10 times. You should get 5 heads. You actually get 8. Is the coin biased, or is it just bad luck?"} Now we have the mathematical tools to solve this without any advanced formulas!
    
    Assume the coin is perfectly fair ($P(H) = 0.5$, $P(T) = 0.5$) and each flip is completely independent.
    \begin{enumerate}[label=\arabic*.]
        \item Imagine one specific sequence of 10 flips: getting 8 heads in a row, followed by 2 tails ($H-H-H-H-H-H-H-H-T-T$). Using the Multiplication Law for independent events, calculate the exact probability of this specific sequence occurring.
        \item Of course, $H-H-H-H-H-H-H-H-T-T$ is not the only way to get 8 heads. You could get $T-H-H-H-H-H-H-H-H-T$, etc. The mathematical formula for finding the number of possible combinations is $\frac{n!}{k!(n-k)!}$. Using $n=10$ and $k=8$, calculate how many different ways you can arrange 8 heads and 2 tails.
        \item Because each of these combinations is \textbf{mutually exclusive}, we can use the Addition Law. Multiply your answer from Part 1 by your answer from Part 2 to find the total probability of getting exactly 8 heads.
        \item Does this percentage prove the coin is mathematically biased, or does it prove that 8 heads can happen purely by chance? 
    \end{enumerate}

\end{enumerate}